\documentclass[a4paper, titlepage]{book}

\usepackage{../preamble}

\begin{document}

\frontmatter

\titlePage{3}{Waves and Oscillations}

\tableofcontents \label{contents}

\mainmatter

\chapter{Springs and Pendulums}

\mathTopics{\pgRef{chap:complex-and-imaginary-numbers}}

\section{Pendulums}

Simple pendulum have a massles rod/string;
therefore, they have no moment of inertia.
By Newton's second law,
\begin{equation}
    \begin{aligned}
        F = m a \\
        -m g \sin \theta = m \ell \ddot{\theta}.
    \end{aligned}
\end{equation}

For angles less than approximately \qty{0.24}{\radian}
(\(\approx\qty{14.3}{\degree}\)),
the small-angle approximation results in a error of less than 1\% in
the gravitational force.
Thus, we find
\begin{align}
    \sin\theta \approx \theta \\
    \therefore - m g \theta            &\approx m \ell \ddot{\theta} \\
    \ddot{\theta} + \omega^2 \theta &\approx 0,~\omega = \sqrt{\frac{g}{\ell}}.
\end{align}

For physical pendulums, the rotational inertia is not 0.
We find
\begin{align}
    -m g d \sin\theta                &=       I \ddot{\theta} \\
    \ddot{\theta} + \omega^2\theta &\approx 0,~\omega = \sqrt{\frac{mgd}{I}},
\end{align}
where \(d\) is the distance between the center of mass and pivot.

\section{Undriven, Undamped Springs}

Simple springs operate are governed by only Newton's second law and Hooke's law.
The equation of motion is
\begin{equation}
    \begin{aligned}
        F              &= m a = m \ddot{x} \\
        F              &= - k x \\
        m \ddot{x}      &= - k x \\
        m \ddot{x} + k x &= 0.
    \end{aligned}
\end{equation}
We can solve this properly using differential equations,
but we are going to skip to the general solutions:
\begin{itemize}
    \item
        \begin{equation}
            x(t) = A \cos\ab(\omega_0 t) + B\sin\ab(\omega_0 t)
        \end{equation},
    \item
        \begin{equation}
            x(t) = A \cos\ab(\omega_0 t + \phi)
        \end{equation},
    \item
        \begin{equation}
            x(t) = z_0 e^{i \omega_0 t} + z^*_0 e^{i \omega_0 t}
        \end{equation},
    \item
        \begin{equation}
            \tilde{x}(t) = A e^{i \phi} e^{i \omega_0 t}
        \end{equation},
\end{itemize}
where
\begin{itemize}
    \item \(\omega_0 = \sqrt{\frac{k}{m}}\) is the \textbf{natural frequency},
    \item \(\Re(\tilde{x}(t)) = x(t)\),
    \item \(A\) and \(B\) are amplitudes,
    \item \(\phi\) is a phase shift,
    \item \(z_0\) is a complex amplitude, which combines both
        amplitude and phase shift.
\end{itemize}
\(A\), \(B\), \(\phi\), and \(z_0\)
can be found by plugging in and solving for initial conditions,
i.e., initial position and velocity.

Solutions to differential equations are \textbf{linear},
i.e., linear combinations of solutions are \emph{also} solutions.
\(\therefore \tilde{x}(t) = A e^{i \phi} e^{i \omega_0 t} + B e^{i
\phi} e^{i \omega_0 t}\)
may also be valid.

Finding velocity and acceleration is straightforward: simply
differentiate position.
We find
\begin{equation}
    \begin{aligned}
        \tilde{v}(t) &= \odv{}{t} \tilde{x}(t) \\
        &= i \omega_0 A e^{i \phi} e^{\omega_0 t} \\
        &= i \omega_0 \tilde{x}(t).
    \end{aligned}
\end{equation}
for velocity.
The factor of \(i\) gives a \qty{90}{\degree} phase shift ahead of position.

Differentiaing again, we find that the acceleration is
\begin{equation}
    \begin{aligned}
        \tilde{a}(t) &= \odv{}{t} \tilde{v}(t) \\
        &= - \omega_0^2 A e^{i \phi} e^{\omega_0 t} \\
        &= - \omega_0^2 \tilde{x}(t).
    \end{aligned}
\end{equation}
The minus indicates a \qty{180}{\degree} phase shift ahead of position.

\section{Undriven, Damped Springs}

Realistically, springs eventually lose energy.
For our purposes, we assume a velocity-dependent spring \textbf{damping}~\(b\)
with units of \unit{\newton \second \per \metre},
and the equation of motion is
\begin{equation}
    \begin{aligned}
        m \ddot{x} = - k x - b \dot{x} \\
        m \ddot{x} + b \dot{x} + k x = 0 \\
        \ddot{x} + \frac{b}{m} \dot{x} + \frac{k}{m} x = 0 \\
        \ddot{x} + \gamma \dot{x} + \omega_0^2 x = 0.
    \end{aligned}
\end{equation}

Some textbooks define \(\gamma\) as \(\frac{b}{2m}\).

We plug a guess of the form \(\boldsymbol{\hat{x}}(t) = A e^{i \phi}
e^{i \beta t}\) into the equation
and get
\begin{equation}
    \boldsymbol{\hat{x}}(t)
    = A
    e^{i \phi}
    e^{- \frac{\gamma}{2} t}
    e^{
        \pm i
        \ab(\sqrt{\omega_0^2 - \frac{\gamma^2}{4}})
        t
    },
\end{equation}
where
\begin{equation}
    \beta = \frac{i \gamma \pm \sqrt{4\omega_0^2 - \omega^2}}{2}.
\end{equation}

Damped oscillators have three possible behaviors
depending on the ratio between the damping coefficient \(\gamma\) and
the natural frequency \(\omega_0\).
The \textbf{quality factor} or \textbf{\(\boldsymbol{Q}\) factor} is defined as
\begin{equation}
    Q \coloneq \frac{\omega_0}{\gamma}.
\end{equation}

The three cases are as follows.

\subsection{Underdamped}

The spring is \textbf{underdamped} when
\begin{equation}
    \begin{aligned}
        \gamma                          &< 2 \omega_0, \\
        \omega_0^2 - \frac{\gamma^2}{4} &> 0, \\
        Q > \frac{1}{2}.
    \end{aligned}
\end{equation}

The position is given by
\begin{equation}
    \tilde{x}(t) = A e^{i \phi} e^{- \frac{\gamma}{2} t} e^{\pm i \omega_u t},
\end{equation}
where
\begin{equation}
    \omega_u = \sqrt{\omega_0^2 - \frac{\gamma^2}{4}}
\end{equation}
is the \textbf{damped frequency}.

\(e^{- \frac{\gamma}{2} t}\) forms an envelope for the solution,
i.e., \(\ab\|x(t)\|\le A e^{- \frac{\gamma}{2} t}\).

Note that although the amplitude decays, the frequency \emph{does not}.

\subsection{Overdamped}

The spring is \textbf{overdamped} when
\begin{equation}
    \begin{aligned}
        \gamma                          &> 2 \omega_0, \\
        \omega_0^2 - \frac{\gamma^2}{4} &< 0, \\
        Q < \frac{1}{2}.
    \end{aligned}
\end{equation}

The position is given by
\begin{equation}
    \boldsymbol{\hat{x}}(t) = A e^{i \phi} e^{- \mu t},
\end{equation}
where
\begin{equation}
    \mu = \frac{\gamma}{2} \pm \sqrt{\frac{\gamma^2}{4} - \omega_0^2}
\end{equation}
is the damped frequency.

\subsection{Critically Damped}

The spring is \textbf{critically damped} when
\begin{equation}
    \begin{aligned}
        \gamma                          &= 2 \omega_0, \\
        \omega_0^2 - \frac{\gamma^2}{4} &= 0, \\
        Q = \frac{1}{2}.
    \end{aligned}
\end{equation}

\section{Driven, Undamped Springs}

\input{chapters/electrical-circuits.tex}
\input{chapters/systems-of-oscillators.tex}
\input{chapters/transverse-waves.tex}
\input{chapters/longitudinal-waves.tex}
\input{chapters/dispersion.tex}
\input{chapters/electromagnetic-waves.tex}

\appendix

\chapter{Complex and Imaginary Numbers}
\label{chap:complex-and-imaginary-numbers}

\(i = \sqrt{-1}\) is the \textbf{imaginary unit}.

\textbf{Imaginary numbers} are the set of all numbers that can be
written in the form \(bi\),
where \(b \in \mathbb{R}\).

\textbf{complex numbers}, denoted by \(\mathbb{C}\),
are the set of all numbers that can be written in the form \(a + bi\),
where \(a,~b \in \mathbb{R}\).

\(a\) is known as the \textbf{real part},
and \(b\) is known as the \textbf{imaginary part}.
Note that this means the imaginary part is \emph{real}.
Given the complex number \(z = a + bi\),
\(\Re(z) = a\), and \(\Im(z) = b\).

The \textbf{complex conjugate} of \(z = a + bi\) is given by \(z^* = a - bi\).

\section{Complex Plane}

Complex numbers can be represented as 2D vectors on a plane,
with the x-axis being the real part (i.e., \(a\)),
and the y-axis being the imaginary part (i.e., \(b\)).

Given a complex number \(z\),
its magnitude \(\ab\| z \|\) is given by Pythagorean's Theorem:
\begin{equation}
    \ab\| z \|
    = \sqrt{\Re(z)^2 + \Im(z)^2}
    = \sqrt{z z^*}
\end{equation}

Its angle \(\phi\) is derived through trigonometry
(be careful about the general limitations of \(\arctan\) on calculators):
\begin{equation}
    \phi = \arctan\ab( \frac{\Im(z)}{\Re(z)} )
\end{equation}

\section{Euler's Formula}

Euler's formula:
\begin{equation}
    e^{i \theta} = \cos\theta + i \sin\theta
\end{equation}

\begin{equation}
    \therefore e^{a + bi} = e^a e^{bi} = e^a \ab( \cos b + i \sin b )
\end{equation}

Oftentimes, working with the complex exponential is easier than the
trigonometric functions.

\section{Complex-Valued Functions}

A \textbf{complex-valued function} is a function that returns a complex number.
It is denoted with a overhead tilde, e.g., \(\tilde{f}(x)\).

\chapter{Fourier Analysis}

\section{Fourier Series}

\section{Fourier Transform}

\backmatter

\printbibliography[
    heading=bibintoc,
    title={References}
]

\end{document}
